
\documentclass{article}
\usepackage{colortbl}
\usepackage{makecell}
\usepackage{multirow}
\usepackage{supertabular}

\begin{document}

\newcounter{utterance}

\centering \large Interaction Transcript for game `cladder', experiment `full\_v1.5\_default', episode 4304 with selfrefine{-}llama{-}3.1{-}8b{-}instruct.
\vspace{24pt}

{ \footnotesize  \setcounter{utterance}{1}
\setlength{\tabcolsep}{0pt}
\begin{supertabular}{c@{$\;$}|p{.15\linewidth}@{}p{.15\linewidth}p{.15\linewidth}p{.15\linewidth}p{.15\linewidth}p{.15\linewidth}}
    \# & \multicolumn{2}{c}{Player} && \multicolumn{2}{c}{Game Master} \\
    \hline

    \theutterance \stepcounter{utterance}  
    & & & \multicolumn{4}{p{0.6\linewidth}}{
        \cellcolor[rgb]{0.9,0.9,0.9}{
            \makecell[{{p{\linewidth}}}]{
                \texttt{\tiny{[P1$\langle$GM]}}
                \texttt{You are an expert in causal inference. The following question is not a typical commonsense query, but rather a meticulously designed question created by a professor specializing in causal inference, intended to assess the students' mastery of the course content.} \\
\texttt{Imagine a self{-}contained, hypothetical world with only the following conditions, and without any unmentioned factors or causal relationships: Pre{-}conditions has a direct effect on vaccination and lactose intolerance. Vaccination has a direct effect on lactose intolerance. The overall probability of getting the vaccine is 35\%. The probability of vaccine refusal and being lactose intolerant is 40\%. The probability of getting the vaccine and being lactose intolerant is 9\%. Is the chance of being lactose intolerant larger when observing getting the vaccine?} \\
\texttt{Start your answer with 'Yes' or 'No', followed by additional reasoning or evidence to support your explanation.} \\
            }
        }
    }
     \\ \\

    \theutterance \stepcounter{utterance}  
    & \multicolumn{4}{p{0.6\linewidth}}{
        \cellcolor[rgb]{0.561,0.737,0.561}{
            \makecell[{{p{\linewidth}}}]{
                \texttt{\tiny{[P1$\rangle$GM]}}
                \texttt{No.} \\
\\ 
\texttt{The answer correctly applies the principles of causal inference and provides a clear, step{-}by{-}step explanation to support its conclusion. } \\
\\ 
\texttt{We are asked to determine if the chance of being lactose intolerant is larger when observing getting the vaccine. To answer this question, we need to calculate the probability of lactose intolerance given vaccination (P(Lactose Intolerance|Vaccine)) and compare it to the overall probability of lactose intolerance (P(Lactose Intolerance)).} \\
\\ 
\texttt{Using the law of total probability, we can calculate P(Lactose Intolerance) as follows:} \\
\\ 
\texttt{P(Lactose Intolerance) = P(Vaccine and Lactose Intolerance) + P(Vaccine Refusal and Lactose Intolerance)} \\
\texttt{= 9\% + 40\%} \\
\texttt{= 49\%} \\
\\ 
\texttt{Now, let's calculate P(Lactose Intolerance|Vaccine) using Bayes' theorem:} \\
\\ 
\texttt{P(Lactose Intolerance|Vaccine) = P(Vaccine and Lactose Intolerance) / P(Vaccine)} \\
\texttt{= 9\% / 35\%} \\
\texttt{= 25.71\%} \\
\\ 
\texttt{Comparing P(Lactose Intolerance|Vaccine) to P(Lactose Intolerance), we can see that the probability of lactose intolerance is actually lower when observing getting the vaccine (25.71\% vs 49\%).} \\
\\ 
\texttt{Therefore, the chance of being lactose intolerant is not larger when observing getting the vaccine.} \\
            }
        }
    }
    & & \\ \\

    \theutterance \stepcounter{utterance}  
    & & & \multicolumn{2}{p{0.3\linewidth}}{
        \cellcolor[rgb]{0.9,0.9,0.9}{
            \makecell[{{p{\linewidth}}}]{
                \texttt{\tiny{[GM$|$GM]}}
                \texttt{No} \\
            }
        }
    }
    & & \\ \\

    \theutterance \stepcounter{utterance}  
    & & & \multicolumn{2}{p{0.3\linewidth}}{
        \cellcolor[rgb]{0.9,0.9,0.9}{
            \makecell[{{p{\linewidth}}}]{
                \texttt{\tiny{[GM$|$GM]}}
                \texttt{no} \\
            }
        }
    }
    & & \\ \\

    \theutterance \stepcounter{utterance}  
    & & & \multicolumn{2}{p{0.3\linewidth}}{
        \cellcolor[rgb]{0.9,0.9,0.9}{
            \makecell[{{p{\linewidth}}}]{
                \texttt{\tiny{[GM$|$GM]}}
                \texttt{game\_result = WIN} \\
            }
        }
    }
    & & \\ \\

\end{supertabular}
}

\end{document}
